\section{Modelagem do Problema}

Nesta secao, a estrutura decisoria introduzida anteriormente e formalizada por
duas representacoes complementares. A primeira descreve o problema completo
como um modelo de particionamento de conjuntos sobre intervalos contiguos. A
segunda apresenta uma versao simplificada em programacao dinamica, util para
explicitar a logica sequencial do problema e para evidenciar com precisao qual
parte da dificuldade decorre da restricao global de disponibilidade da frota.

\subsection{Formulacao por particionamento de conjuntos sobre intervalos}

Uma forma natural de representar o problema e tratar cada decisao viavel de
carregamento como uma viagem candidata. Cada viagem corresponde a um bloco
contiguo de cargas atendido por um veiculo de determinada classe. A solucao
consiste, entao, em selecionar um conjunto de viagens que particione toda a
sequencia fisica disponivel, respeite a disponibilidade da frota e minimize o
custo total da operacao.

Seja \(N = \{1,\dots,n\}\) o conjunto das cargas, na ordem em que estao
fisicamente disponiveis, e seja \(K = \{1,\dots,m\}\) o conjunto das classes de
veiculos. Para cada carga \(t \in N\), definem-se seu peso \(w_t\) e seu
comprimento \(\ell_t\). Para cada classe \(k \in K\), definem-se a capacidade de
peso \(W_k\), a capacidade de comprimento \(L_k\), a quantidade disponivel
\(M_k\) e o frete minimo faturavel \(q_k\).

Para cada bloco contiguo \(i,\dots,j\), com \(1 \le i \le j \le n\), definem-se
os acumulados
\[
\bar{W}_{ij} = \sum_{t=i}^{j} w_t
\qquad \text{e} \qquad
\bar{L}_{ij} = \sum_{t=i}^{j} \ell_t.
\]
Se o bloco \(i,\dots,j\) for compativel com um veiculo da classe \(k\), define-se
o custo da viagem candidata como
\[
c_{ijk} = \max(\bar{W}_{ij}, q_k),
\]
capturando o fato de que a cobranca depende do peso transportado, mas pode ser
limitada inferiormente por um valor minimo contratual.

Introduz-se a variavel binaria \(x_{ijk}\), igual a 1 quando a classe de
veiculo \(k\) e usada para transportar o bloco contiguo \(i,\dots,j\), e igual a
0 caso contrario. Considerando apenas combinacoes viaveis, uma formulacao
inicial do problema pode ser escrita como:

\begin{align}
\min \quad & \sum_{(i,j,k)} c_{ijk}\, x_{ijk} \\
\text{sujeito a} \quad
& \sum_{(i,j,k):\, i \le t \le j} x_{ijk} = 1
\quad \forall t \in N \\
& \sum_{(i,j)} x_{ijk} \le M_k
\quad \forall k \in K \\
& x_{ijk} \in \{0,1\}.
\end{align}

A primeira restricao garante que cada carga apareca exatamente uma vez na
solucao, o que assegura simultaneamente cobertura total da sequencia e ausencia
de sobreposicoes. A segunda limita o numero de viagens selecionadas por classe
de veiculo conforme a disponibilidade da frota. Conceitualmente, trata-se de uma
formulacao do tipo \emph{set partitioning over intervals}, em que cada coluna
representa um intervalo contiguo viavel associado a uma classe de veiculo.

\subsection{Versao simplificada em programacao dinamica}

A mesma estrutura tambem pode ser lida de forma recursiva. Em vez de selecionar
todas as viagens de uma unica vez, constroi-se a solucao da esquerda para a
direita, sempre fechando a ultima viagem de um prefixo da sequencia. Essa
interpretacao preserva a logica de particionamento em blocos contiguos, mas
torna mais transparente a relacao entre a formulacao inteira e uma abordagem em
programacao dinamica.

Defina \(dp(j)\) como o melhor custo para cobrir exatamente as cargas
\(1,\dots,j\). O caso base e dado por
\[
dp(0) = 0.
\]
Para cada posicao \(j\), a ultima viagem da solucao otima para o prefixo
\(1,\dots,j\) deve necessariamente corresponder a algum bloco contiguo
\(i,\dots,j\) servido por uma classe de veiculo viavel \(k\). Isso leva a
recorrencia
\[
dp(j) =
\min_{\substack{1 \le i \le j \\ k \in K \text{ viavel}}}
\left\{ dp(i-1) + c_{ijk} \right\}.
\]
Mantendo-se, para cada \(j\), o par \((i,k)\) que realiza o minimo, a solucao
completa pode ser reconstruida por retrocesso a partir de \(dp(n)\).

Essa recorrencia mostra de forma explicita a conexao com a formulacao anterior.
No modelo de particionamento, escolhe-se um subconjunto de viagens candidatas
que cobre toda a sequencia. Na programacao dinamica simplificada, escolhe-se a
ultima viagem de cada prefixo e encadeiam-se essas decisoes ate o final. Em
ambos os casos, os objetos fundamentais sao os mesmos: intervalos contiguos
viaveis e seus custos associados.

Essa leitura tambem admite uma interpretacao em grafo aciclico. Considerem-se
os nos \(0,1,\dots,n\), onde cada arco de \(i-1\) para \(j\) representa o bloco
contiguo \(i,\dots,j\). O custo do arco corresponde ao custo da viagem
candidata escolhida para esse bloco. Sob essa perspectiva, a recorrencia acima
equivale a um problema de caminho minimo em rede aciclica.

\subsection{Limites da versao simplificada e relacao entre os modelos}

Apesar de sua clareza estrutural, a programacao dinamica descrita acima nao e
equivalente ao problema completo quando a disponibilidade de frota e uma
restricao efetiva. O motivo e direto: o estado \(dp(j)\) registra apenas ate que
posicao da sequencia se chegou e o melhor custo correspondente, mas nao informa
quantos veiculos de cada classe ja foram consumidos ao longo da solucao parcial.

Essa omissao e decisiva. Duas solucoes parciais podem cobrir exatamente o mesmo
prefixo \(1,\dots,j\) com custos semelhantes, mas utilizar a frota de maneiras
distintas. Se uma delas consome prematuramente uma classe de veiculo escassa,
ela pode inviabilizar ou encarecer as decisoes futuras, ainda que pareca
atraente do ponto de vista local. Como o estado simplificado nao distingue essas
situacoes, a recorrencia trata como equivalentes solucoes parciais que, no
problema completo, nao o sao.

Consequentemente, a versao simplificada em programacao dinamica resolve de forma
exata apenas uma relaxacao do modelo inteiro: o caso em que a restricao de frota
nao existe, e irrelevante, ou nunca se torna ativa. Nessa leitura, a dificuldade
adicional do problema completo nao decorre da contiguidade nem do calculo dos
custos de cada bloco, mas da interacao entre decisoes locais de particionamento
e o consumo global de recursos escassos.

Essa comparacao e util por duas razoes. Primeiro, ela ajuda a organizar a
estrutura teorica do problema: o modelo inteiro e a formulacao mais fiel ao
contexto operacional completo, enquanto a programacao dinamica simplificada
isola a logica sequencial central do problema. Segundo, ela sugere caminhos para
o desenvolvimento futuro da pesquisa, seja por meio de estados enriquecidos em
programacao dinamica, seja por reformulacoes em rede, decomposicao ou uso da
versao simplificada como baseline computacional, heuristica ou limite inferior.
